\section{Folio 2}

\begin{ejercicio}
    Para el operador $T: \ell_p \longrightarrow \ell_p$, con $1 \leq p < \infty$, definido por \\
    $T(x) = \left( x_1, \frac{x_2}{2}, \dots, \frac{x_n}{n}, \dots \right)$ para $x = (x_1, \dots, x_n, \dots) \in \ell_p$, demostrar:
    \begin{itemize}
        \item[(i)] $0 \in \sigma(T)$
        \item[(ii)] $0$ no es un valor propio
        \item[(iii)] $\sigma(T) = \left\{ 0, 1, \dots, \frac{1}{n}, \dots \right\}$
    \end{itemize}
    \begin{description}
        \item[$i)$] Para demostrar que $0 \in \sigma(T)$, para ello sabemos primero que $\ell_p$ es de dimensión infinita, por lo tanto si vemos que el operador es compacto tendríamos lo pedido. Para ello, basta con observar que $T$ es el límite de los operadores de rango finito $T_n: \ell_p \longrightarrow \ell_p$ (que son compactos) definidos por:
        $$T_k(x) = \left( x_1, \frac{x_2}{2}, \dots, \frac{x_k}{k}, 0, 0, \dots \right)$$
        Ahora calculamos la norma de la diferencia $(T - T_k)$
        $$(T - T_k)(x) = \left( 0, \dots, 0, \frac{x_{k+1}}{k+1}, \frac{x_{k+2}}{k+2}, \dots \right)$$
        Evaluamos su norma en $\ell_p$
        $$\|(T - T_k)(x)\|_p^p = \sum_{n=k+1}^{\infty} \left| \frac{x_n}{n} \right|^p \leq \frac{1}{(k+1)^p} \sum_{n=k+1}^{\infty} |x_n|^p \leq \frac{1}{(k+1)^p} \|x\|_p^p$$
        Tomando la raíz $p$-ésima a ambos lados
        $$\|(T - T_k)(x)\|_p \leq \frac{1}{k+1} \|x\|_p$$
        Por la definición de norma de un operador
        $$\|T - T_k\| = \sup_{x \neq 0} \frac{\|(T - T_k)(x)\|_p}{\|x\|_p} \leq \frac{1}{k+1}$$
        Si hacemos el límite cuando $k \to \infty$
        $$\lim_{k \to \infty} \|T - T_k\| = 0$$
        Como $T$ es el límite en norma de una sucesión de operadores compactos ($T_k$), $T$ es compacto, y se tendría $0\in \sigma(T)$.

        \item[$ii)$] La idea a aplicar es que cuando decimos que $\lambda = 0$ es un valor propio significa que existe un vector no nulo $x \neq 0$ tal que $T(x) = 0 \cdot x = 0$. Es decir, buscamos si el núcleo $\ker(T)$ tiene más elementos que el cero. Veámoslo rigurosamente. \\
        \noindent
        Planteamos la ecuación de autovalores para $\lambda = 0$:$$T(x) = 0 \implies \left( x_1, \frac{x_2}{2}, \dots, \frac{x_n}{n}, \dots \right) = (0, 0, \dots, 0, \dots)$$
        Igualando componente a componente:  $x_1 = 0$$\frac{x_2}{2} = 0 \implies x_2 = 0$$\dots$$\frac{x_n}{n} = 0 \implies x_n = 0 \quad \forall n \in \mathbb{N}$. El único vector que cumple esto es el vector nulo $x = 0$. Como no existe ningún vector $x \neq 0$ tal que $T(x) = 0$, concluimos rigurosamente que $0$ no es un valor propio ($0 \notin e(T)$). 
        \item[$iii)$] Sea $\lambda \in \mathbb{C}$ tal que $\lambda \neq 0$ y $\lambda \neq \frac{1}{n}$ para todo $n \in \mathbb{N}$. Queremos probar que $\lambda \in \rho(T)$, es decir, que para cualquier sucesión $y \in \ell_p$, existe una única solución $x \in \ell_p$ tal que
        $$(T - \lambda I)(x) = y$$
        donde lo que estamos demostrando a priori es la sobreyectividad e inyectividad. Si escribimos esta ecuación componente a componente:
        \begin{equation*}
            \left( \left(1 - \lambda\right)x_1, \left(\frac{1}{2} - \lambda\right)x_2, \dots, \left(\frac{1}{n} - \lambda\right)x_n, \dots \right) = (y_1, y_2, \dots, y_n, \dots)
        \end{equation*}
        Como hemos elegido un $\lambda$ que no anula ningún término (ya que $\lambda \neq \frac{1}{n}$), podemos despejar cada componente $x_n$ de forma única
        $$x_n = \frac{y_n}{\frac{1}{n} - \lambda} \quad \forall n \in \mathbb{N}$$
        falta ver que efectivamente $(x_n)\in \ell_p$. Para conseguirlo, analizamos el comportamiento del denominador cuando $n \to \infty$
        $$\lim_{n \to \infty} \left( \frac{1}{n} - \lambda \right) = 0 - \lambda = -\lambda$$Como $\lambda \neq 0$, este límite es un número estrictamente distinto de cero. Esto significa que la sucesión de los denominadores no se acerca a cero infinitamente; se mantiene ``lejos'' de cero a partir de cierto término. Por lo tanto, la sucesión inversa está acotada por una constante superior $M > 0$
        $$\left| \frac{1}{\frac{1}{n} - \lambda} \right| \leq M \quad \forall n \in \mathbb{N}$$
        Por lo tanto, se cumple la siguiente desigualdad
        $$\sum_{n=1}^{\infty} \left| \frac{y_n}{\frac{1}{n} - \lambda} \right|^p \leqq M^p \sum_{n=1}^{\infty} |y_n|^p < +\infty$$
        lo que implica que $x = (x_n) \in \ell_p$. 
    \end{description}
\end{ejercicio}

\begin{ejercicio}
    Sea $H$ un espacio de Hilbert y $M$ un subespacio cerrado de $H$, con $M \neq H$. Demostrar que si $P = P_M$ es el operador de proyección ortogonal sobre $M$, entonces los números $0$ y $1$ son valores propios de $P$. \\
    Demostrar además que el conjunto de los valores propios es precisamente $\{0,1\}$ \big(en realidad se podría demostrar que $\sigma(P) = \{0,1\}$\big). \\

    \noindent
    Usaremos el Teorema de la Proyección Ortogonal (o de Descomposición Ortogonal), que dice que si $M$ es un subespacio cerrado de un espacio de Hilbert $H$, entonces el espacio entero se puede romper limpiamente en la suma directa de $M$ y su complemento ortogonal $M^\perp$, es decir
    $$H = M \oplus M^\perp$$
    lo que significa que cualquier vector $x \in H$ se puede escribir de forma única como
    $$x = x_M + x_{M^\perp}$$
    donde $x_M \in M$ y $x_{M^\perp} \in M^\perp$. El operador de proyección ortogonal actúa quedándose únicamente con la parte que vive en $M$: $P(x) = x_M$. \\  \\
    \noindent
    Demostraremos ahora que $0$ y $1$ son valores propios de $P$. Para probar que un número $\lambda$ es un valor propio, debemos encontrar un vector no nulo $x \neq 0$ tal que $P(x) = \lambda x$.
    \begin{itemize}
        \item \tbf{Caso $\lambda=1$}: tomemos un vector cualquiera $x$ que viva dentro del subespacio $M$ (con $x \neq 0$, lo cual es posible porque al ser subespacio contiene al menos al vector nulo, y el enunciado no dice que $M = \{0\}$). \\
        \noindent
        Como $x \in M$, su descomposición única es simplemente $x = x + 0$, donde su componente ortogonal es cero. Por definición del operador de proyección:$$P(x) = x = 1 \cdot x$$Por lo tanto, $\lambda = 1$ es un valor propio de $P$, y cualquier vector de $M$ es un autovector asociado (es decir, $M \subseteq V_1$).
        \item \tbf{Caso $\lambda = 0$}: tomemos ahora un vector cualquiera $y$ que viva en el complemento ortogonal, $y \in M^\perp$. Como el enunciado nos dice que $M \neq H$, el complemento ortogonal $M^\perp$ no es vacío y contiene vectores distintos de cero ($y \neq 0$). \\
        \noindent
        La descomposición única de este vector es $y = 0 + y$, ya que su componente en $M$ es nula. Aplicando el operador de proyección:$$P(y) = 0 = 0 \cdot y$$Por lo tanto, $\lambda = 0$ es un valor propio de $P$, y cualquier vector de $M^\perp$ es un autovector asociado (es decir, $M^\perp \subseteq V_0$)
    \end{itemize}
    Queremos probar que no puede existir ningún otro valor propio misterioso fuera de estos dos. Para ello, supongamos que existe un valor propio $\lambda \in \mathbb{C}$ tal que $\lambda\neq 0,1$ con un vector propio asociado no nulo $x \neq 0$, tal que
    $$P(x) = \lambda x$$
    Vamos a descomponer de forma única a ese autovector $x$ usando el teorema de la proyección: $x = x_M + x_{M^\perp}$, por lo que si aplicamos el operador $P$ a ambos lados de la igualdad de autovalores, obtenemos por un lado
    $$P(x) = P(x_M + x_{M^\perp}) = P(x_M) + P(x_{M^\perp}) = x_M + 0 = x_M$$
    Y por el otro lado, sustituyendo $x$ por su descomposición en la expresión $\lambda x$
    $$\lambda x = \lambda (x_M + x_{M^\perp}) = \lambda x_M + \lambda x_{M^\perp}$$
    Igualando ambas expresiones ($P(x) = \lambda x$)
    $$x_M = \lambda x_M + \lambda x_{M^\perp} \implies (1 - \lambda)x_M - \lambda x_{M^\perp} = 0$$
    Como $M$ y $M^\perp$ son ortogonales, la intersección entre ellos es únicamente el vector cero (\mbox{$M \cap M^\perp = \{0\}$}). Dado que $(1 - \lambda)x_M \in M$ y $\lambda x_{M^\perp} \in M^\perp$, para que su resta sea el vector nulo, cada componente por separado debe ser obligatoriamente cero
    \begin{equation*}
        (1 - \lambda)x_M = 0 \qquad \text{y} \qquad \lambda x_{M^\perp} = 0
    \end{equation*}
    Como el autovector original $x \neq 0$, es imposible que $x_M$ y $x_{M^\perp}$ sean cero a la vez. Esto nos fuerza a analizar dos únicos caminos posibles:
    \begin{itemize}
        \item Si $x_M \neq 0$, entonces de la ecuación $(1 - \lambda)x_M = 0$ se deduce que $(1 - \lambda) = 0$, lo que implica que $\lambda = 1$, contradicción.
        \item Si $x_{M^\perp} \neq 0$, entonces de la ecuación $\lambda x_{M^\perp} = 0$ se deduce que $\lambda = 0$, contradicción.
    \end{itemize}
    por lo que no puede existir ningún valor propio $\lambda$ distinto de $0$ y $1$. Por lo tanto, el conjunto de los valores propios de $P$ es exactamente $\{0,1\}$, es decir
    \begin{equation*}
        e(T)=\{0,1\}
    \end{equation*}
\end{ejercicio}

\begin{ejercicio} 
    Sea $H$ un espacio de Hilbert y $T: H \longrightarrow H$ un operador compacto. Demostrar que no puede existir una constante $\alpha > 0$ tal que:
    \[ ( T(x), x ) \geq \alpha \|x\|^2, \qquad \forall x \in H \]

    \noindent
    Es importante notar que vamos a suponer dimensión infinita para $\mathcal{H}$, pues para el caso de \underline{dimensión finita}, no es cierto lo propuesto. Podemos tomar por ejemplo el espacio de Hilbert $\R^2$ de dimensión finita, y donde sabemos que la bola unidad es compacta y entonces 

    \Func{Id}{\R^2}{\R^2}{x}{x}
    es compacto. Estaríamos en la hipótesis del enunciado, pero en cambio tendríamos
    $$ ( Tx, x ) = ( x, x ) =1\cdot \|x\|^2$$
    y tomando $\alpha=1$ tendríamos la contradicción del enunciado. \\

    \noindent
    Dicho esto, estudiemos ahora el caso de $\mathcal{H}$ \underline{dimensión infinita}, donde supondremos que si existe $\alpha>0$ que cumple dicha característica y veremos a donde llegamos.
    
    \noindent
    Como el espacio de Hilbert $H$ es de dimensión infinita, existe un sistema ortonormal infinito de vectores $\{e_n\}_{n \in \mathbb{N}} \subset H$, que por definición cumplen
    $$\|e_n\| = 1 \quad \forall n \in \mathbb{N} \quad \text{y} \quad \langle e_n, e_m \rangle = 0 \quad \text{si } n \neq m$$
    y de los cuales la distancia entre dos de ellos distintos (ya calculada en teoría varias veces) se puede obtener así
    \begin{align*}
        \|e_n - e_m\|^2 =& ( e_n, e_n ) - ( e_n, e_m ) - ( e_m, e_n ) + ( e_m, e_m )\|e_n - e_m\|^2 = \\
        =& 1 - 2(0) + 1 = 2\implies \|e_n - e_m\| = \sqrt{2} \quad \forall n \neq m
    \end{align*}
    Por propiedades por ser sistema ortonromal, sabemos que la sucesión $\{e_n\}$ está acotada y como $T$ es compacto, podemos obtener una subsucesión de sus imágenes convergente, es decir
    
    $$T(e_{n_k}) \longrightarrow z \quad \text{en norma de } H \text{ cuando } k \to \infty$$   
    \noindent
    Dado que toda sucesión convergente es una sucesión de Cauchy ($\mathcal{H}$ espacio de Hilbert), esta subsucesión de imágenes cumple obligatoriamente que la distancia entre sus términos tiende a cero cuando los índices crecen
    $$\|T(e_{n_k}) - T(e_{n_m})\| \longrightarrow 0 \quad \text{cuando } k, m \to \infty$$
    Vamos a evaluar la propieda de la hipótesis con $x = e_{n_k} - e_{n_m}$ donde $k \neq m$

    \begin{align*}
        \alpha \|e_{n_k} - e_{n_m}\|^2\leq (T(e_{n_k}) - T(e_{n_m}), e_{n_k} - e_{n_m} ) \overset{(*)}{\leq} \|T(e_{n_k}) - T(e_{n_m})\| \cdot \|e_{n_k} - e_{n_m}\|
    \end{align*}
    donde en $(*)$ hemos usado la desigualdad de Cauchy-Schwarz. Aplicando la distancia entre dos elementos de la sucesión distintos tenemos que
    $$\alpha \cdot 2 \leq \|T(e_{n_k}) - T(e_{n_m})\| \cdot \sqrt{2}\implies \sqrt{2} \alpha \leq \|T(e_{n_k}) - T(e_{n_m})\|$$
    donde finalmente si tendemos $k,m\to \infty$ tenemos que
    $$\sqrt{2} \alpha \leq 0 \implies \alpha \leq 0$$
    lo cual es una contradicción, pues dijimos $\alpha>0$, entonces no existe dicho $\alpha$ y queda demostrado lo que se quería.
\end{ejercicio}

\begin{ejercicio}
    Sea $\mathcal{H} = \ell_2$, entonces dado el siguiente operador
    \Func{T}{\mathcal{H}}{\mathcal{H}}{(x_1, x_2, \dots, x_n, \dots)}{(x_2, \dots, x_{n+1}, \dots)}
    Calcular el operador adjunto $T^*$ y demostrar que ningún número complejo $\lambda \in \mathbb{C}$ puede ser valor propio de $T^*$. \\

    \noindent
    Recordemos que el espacio $\mathcal{H} = \ell_2$ es el espacio de sucesiones complejas de cuadrado sumable, dotado del producto escalar estándar
    $$( x, y ) = \sum_{n=1}^{\infty} x_n \overline{y_n}$$
    Por definición, el operador adjunto $T^*$ es el único operador que satisface la identidad
    $$( Tx, y ) = ( x, T^*y ) \quad \forall x, y \in \ell_2$$
    Desarrollemos el lado izquierdo aplicando la definición de tu operador $T(x) = (x_2, x_3, \dots, x_{n+1}, \dots)$
    \begin{align*}
        ( Tx, y ) &= \sum_{n=1}^{\infty} (Tx)_n \overline{y_n} = x_2\overline{y_1} + x_3\overline{y_2} + x_4\overline{y_3} + \dots + x_{n+1}\overline{y_n} + \dots= \\
        &=x_1 \cdot \overline{0} + x_2 \cdot \overline{y_1} + x_3 \cdot \overline{y_2} + \dots + x_n \cdot \overline{y_{n-1}} + \dots =( (x_1, x_2, x_3, \dots), (0, y_1, y_2, y_3, \dots) )
    \end{align*}
    Comparando con la definición del adjunto $\langle x, T^*y \rangle$, deducimos por unicidad que
    $$T^*(y) = (0, y_1, y_2, \dots, y_{n-1}, \dots)$$
    a este operador $T^*$ se le conoce como el operador de desplazamiento a la derecha (\tbf{right shift operator}).\\ 

    \noindent
    Veremos ahora si hay \underline{valores propios}. Por definición, un número complejo $\lambda \in \mathbb{C}$ es un valor propio de $T^*$ si existe un vector no nulo $y \in \ell_2$ ($y \neq 0$) tal que se cumple la ecuación de autovalores
    $$T^*(y) = \lambda y$$
    donde si la vemos componente a componente viene dada como
    \begin{equation*}
        \begin{cases}
            0&=\lambda y_1 \\
            y_1&=\lambda y_2 \\
            \vdots \\
            y_n&=\lambda y_{n+1} \\
            \vdots
        \end{cases}
    \end{equation*}
    Veámos los distintos casos de $\lambda$:
    \begin{itemize}
        \item \tbf{Caso $\lambda=0$}: de la primera ecuación $0 = \lambda y_1$ se deduce que $y_1$ puede ser cualquier número complejo, pero de la segunda ecuación $y_1 = \lambda y_2$ se deduce que $y_1$ debe ser necesariamente cero. Si seguimos con el resto de componentes vemos que todos los términos de la sucesión sean cero 
        $$y = (0, 0, 0, \dots)$$ 
        y como el vector nulo no puede ser un vector propio por definición, concluimos que $\lambda = 0$ no es valor propio.
        \item \tbf{Caso $\lambda \neq 0$}: de la primera ecuación $0 = \lambda y_1$ se deduce que $y_1 = 0$. De la segunda ecuación $y_1 = \lambda y_2$ se deduce que $y_2 = 0$. De la tercera ecuación $y_2 = \lambda y_3$ se deduce que $y_3 = 0$. Siguiendo este razonamiento, vemos que todos los términos de la sucesión deben ser cero
        $$y = (0, 0, 0, \dots)$$
        y como el vector nulo no puede ser un vector propio por definición, concluimos que ningún $\lambda \neq 0$ es valor propio.
    \end{itemize}
    Finalmente obtenemos que ningún número complejo $\lambda \in \mathbb{C}$ es valor propio de $T^*$, es decir
    $$e(T^*) = \emptyset$$
\end{ejercicio}


\begin{ejercicio}
    Si $X$ es espacio de Banach de dimensión infinita y $T \in K(X)$, entonces 
    $$0 \in \sigma(T)$$
    Si por el contrario (P.A.) $T - 0 \cdot Id = T$ fuera biyectiva, entonces $T \circ T^{-1} = Id \in K(X)$ (ya que la composición de operadores compactos es compacta), veámos que la identidad $Id$ no es compacto.\\

    \noindent
    Sabemos que como $X$ es de dimensión infinita, entonces la bola unitaria no es compacta. Si el operador $Id: X \to X$ fuese compacto, por definición, transformaría el conjunto acotado $B_X$ en un conjunto relativamente compacto. Como $Id(B_X) = B_X$, esto significaría que la clausura de $B_X$ (que es ella misma, por ser cerrada) es compacta. Es decir, $B_X$ sería compacta en norma, pero dijimos que no era compacta, y por tanto como llegamos a una contradicción, tenemos que $Id$ no es compacto, es decir

    $$0 \in \sigma(T)$$
\end{ejercicio}


\begin{ejercicio} % //TODO: hacer
    Sea $T: \ell_2 \longrightarrow \ell_2$ el operador definido por:
    \begin{itemize}
        \item[(I)] $T(x) = \{\alpha_i x_i\}_{i \in \mathbb{N}^+}$ \quad si $x = (x_1, \dots, x_j, \dots) \in \ell_2$ \\
        donde $\{\alpha_j\}$ es una sucesión que representa los números racionales en $[0,1]$, es decir 
        $$\{\alpha_j\} = \mathbb{Q} \cap [0,1]$$
        
        \item[(II)] $T(x) = \left( 0, x_1, \frac{x_2}{2}, \dots, \frac{x_k}{k}, \dots \right)$ \quad si $x = (x_1, \dots, x_j, \dots) \in \ell_2$
    \end{itemize}
    En ambos casos, I) o II), determinar el espectro de $T$ y los posibles valores propios. ¿Es $T$ compacto? \\

    \noindent
    Veámos ambos casos por separado.
    \begin{description}
        \item[Caso I)] Como $\{\alpha_j\}$ es una numeración de los racionales en $[0,1]$, todos los coeficientes cumplen que $|\alpha_j| \le 1$, entonces
        \begin{equation*}
            \|T(x)\|_{\ell_2}^2 = \sum_{i=1}^{+\infty} |\alpha_i x_i|^2 = \sum_{i=1}^{+\infty} |\alpha_i|^2 |x_k|^2 \le \sum_{i=1}^{+\infty} 1 \cdot |x_i|^2 = \sum_{i=1}^{+\infty} |x_i|^2=\|x\|_{\ell_2}^2
        \end{equation*}
        que cumple la continuidad de $T$ y la linealidad es trivial. Vamos a probar que $T$ no es compacto, para ellos cogeremos la base canónica $\{e_n\}$, y tenemos que $T(e_n) = \alpha_n e_n$, ahora tomemos una subsucesión cualquier de esta sucesión de imágenes, es decir, $T(e_{n_k}) = \alpha_{n_k} e_{n_k}$, y veamos su comportamiento. 
        
        Dado que $\{\alpha_j\}$ es una numeración de los racionales en $[0,1]$, esta sucesión de coeficientes $\{\alpha_{n_k}\}$ no puede tener ningún límite (pues los racionales no tienen ningún punto de acumulación), por lo que la distancia entre dos términos distintos de esta sucesión no puede tender a cero, es decir
    \end{description}
\end{ejercicio}

\begin{ejercicio}
    Sea el operador
    \Func{T}{C^0([0,1])}{C^0([0,1])}{x}{T(x)}
    donde $x: [0,1] \to \mathbb{R}\in C^0([0,1])$ (función continua), y $T(x)(t) := t\cdot x(t)$. Se pide:
    \begin{itemize}
        \item[i)] Calcular $\sigma(T)$.
        \item[ii)] Demuestra si $T$ es compacto o no.
    \end{itemize}
    Veámos la resolución de cada apartado por separado.
    \begin{description}
        \item[$i)$] Tomemos una función arbitraria $f \in C^0([0,1])$ y un número $\lambda \in \mathbb{R}$, entonces la ecuación de resolvente es la siguiente
        \begin{equation*}            
            (T - \lambda Id)(x) = f
        \end{equation*}
        y buscamos encontrar si existe dicho $x$ que cumpla esta ecuación. Si existe, entonces $\lambda \in \rho(T)$, y si no existe, entonces $\lambda \in \sigma(T)$. Para resolver esta ecuación, basta con pasar todo a un solo lado y sacar factor común $x(t)$, quedando la ecuación de la siguiente forma
        \begin{equation*}
            (T - \lambda Id)(x)(t) = t \cdot x(t) - \lambda x(t) = (t - \lambda) x(t) = f(t) \iff x(t) = \frac{f(t)}{t - \lambda} \quad \forall t \in [0,1]
        \end{equation*}
        Ahora tenemos dos casos a estudiar
        \begin{itemize}
            \item Si $\lambda \notin [0,1]$, entonces $t - \lambda \neq 0$ para todo $t \in [0,1]$, por lo que $x(t) = \frac{f(t)}{t - \lambda}$ es una función continua (ya que $f$ es continua y el denominador no se anula), por lo que existe solución para cualquier $f$, y además es única, por lo que $\lambda \in \rho(T)$.
            \item Si $\lambda \in [0,1]$, entonces existirá un punto exacto en el dominio, al que llamamos $t_0 = \lambda$, donde el denominador se desvanece por completo: $(t_0 - \lambda) = 0$. Si intentamos evaluar la solución candidata en ese punto, tendríamos:$$x(\lambda) = \frac{f(\lambda)}{0}$$Si elegimos una función $f$ tal que $f(\lambda) \neq 0$, la división por cero genera una asíntota vertical (una discontinuidad que tiende a infinito). Es imposible encontrar una función continua $x(t)$ que satisfaga la igualdad en ese punto. Por lo tanto, no existe solución para esa función $f$, y por tanto $\lambda \in \sigma(T)$. Dado esto tenemos entonces que 
            \begin{equation*}
                [0,1]\subseteq \sigma(T)
            \end{equation*}
        \end{itemize}
        Juntando ambos cosas tenemos finalmente que 
        \begin{equation*}
            \sigma(T) = [0,1]
        \end{equation*}
        
        \item[$ii)$] Para determinar si $T$ es compacto o no, basta con estudiar si el espectro de $T$ sin el punto $0$ coincide con el conjunto de valores propios sin el punto $0$. En este caso vamos a demostrar que el conjunto de valores propios es vacío, es decir, $e(T) = \emptyset$, lo cual implica que $T$ no es compacto (de lo contrario se tendría $\sigma(T) \setminus \{0\} = e(T) \setminus \{0\}$). Para ello, basta con resolver la ecuación de valores propios 
        $$T(x) = \lambda x \iff T(x)(t) = \lambda x(t) \iff  t\cdot x(t) = \lambda x(t)\iff (t - \lambda) x(t) = 0 \quad \forall t\in [0,1]$$
        donde tenemos que $x$ debe ser una función no nula. Para que el producto de dos números reales sea cero, es obligatorio que el primero sea cero, el segundo sea cero, o ambos. Por lo tanto, en cada punto $t \in [0,1]$, la matemática nos obliga a elegir entre una de estas dos opciones:
        \begin{itemize}
            \item Opción A: el primer paréntesis se anula: $t = \lambda$.
            \item Opción B: la función se anula en ese instante: $x(t) = 0$.
        \end{itemize}
        Supongamos por un momento que sí existe un valor propio $\lambda$ con su correspondiente función propia $x(t) \neq 0$. Como hemos supuesto que $x \neq 0$, la definición nos garantiza que existe al menos un punto concreto, llamémoslo $t_0$, donde la función tiene una altura distinta de cero
        $$x(t_0) \neq 0$$
        Si en el punto $t_0$ la función se encuentra a una altura $x(t_0) \neq 0$, está obligada a mantenerse alejada del cero en todo un pequeño entorno (un intervalo abierto) alrededor de $t_0$, entonces
        $$x(t) \neq 0 \quad \forall t \in \text{Entorno}(t_0)$$
        lo cual es claramente un absurdo, pues $\lambda$ es una constante y no puede tener varios valores en un entorno de $t_0$. Por lo tanto, no existe ningún valor propio $\lambda$ y $e(T) = \emptyset$.
    \end{description}
\end{ejercicio}

\begin{ejercicio}
    Sea $T \in BL(H)$ autoadjunto y compacto. Entonces los números:
    \[ m = \inf_{\substack{\|u\|=1 \\ u \in H}} \langle T(u), u \rangle \qquad M = \sup_{\substack{u \in H \\ \|u\|=1}} \langle T(u), u \rangle \]
    son mínimo y máximo. \\

    \noindent
    Queremos demostrar que el supremo $M = \sup\limits_{\|u\|=1} ( Tu, u )$ se alcanza, es decir, que existe un vector $\bar{u} \in H$ con $\|\bar{u}\|=1$ tal que $( T\bar{u}, \bar{u} ) = M$. \\

    \noindent
    Por la propia definición de supremo, existe una sucesión de vectores $\{u_n\}_{n \in \mathbb{N}} \subset H$ en la esfera unidad ($\|u_n\| = 1$) tal que sus imágenes por la forma cuadrática convergen a $M$
    $$\lim_{n \to \infty} \langle Tu_n, u_n \rangle = M$$
    Como la sucesión $\{u_n\}$ está acotada ($\|u_n\|=1$), por las propiedades de los espacios de Hilbert, sabemos que posee una subsucesión $\{u_{n_k}\}$ que converge débilmente a algún elemento $\bar{u} \in H$ (es decir, $u_{n_k} \xrightarrow{w} \bar{u}$). \\
    
    \noindent
    Aquí es donde entra la compacidad de $T$. Por definición, un operador es compacto si transforma sucesiones débilmente convergentes en sucesiones fuertemente convergentes (en norma). Por lo tanto
    $$u_{n_k} \xrightarrow{w} \bar{u} \implies T(u_{n_k}) \xrightarrow{\|\cdot\|} T(\bar{u})$$
    Vamos a calcular el límite de nuestra forma cuadrática a lo largo de esta subsucesión
    $$\langle T(u_{n_k}), u_{n_k} \rangle = \langle T(u_{n_k}) - T(\bar{u}) + T(\bar{u}), \; u_{n_k} \rangle = \langle T(u_{n_k}) - T(\bar{u}), u_{n_k} \rangle + \langle T(\bar{u}), u_{n_k} \rangle$$
    Analicemos los dos términos cuando $k \to \infty$:
    \begin{itemize}
        \item El primer término se apaga por la convergencia fuerte en norma del Paso 2
        $$|\langle T(u_{n_k}) - T(\bar{u}), u_{n_k} \rangle| \le \|T(u_{n_k}) - T(\bar{u})\| \cdot \|u_{n_k}\| \longrightarrow 0$$
        \item El segundo término converge por la definición de convergencia débil ($u_{n_k} \xrightarrow{w} \bar{u}$)
        $$\lim_{k \to \infty} \langle T(\bar{u}), u_{n_k} \rangle = \langle T(\bar{u}), \bar{u} \rangle$$
    \end{itemize}
    Uniendo ambos resultados, obtenemos que
    $$M = \lim_{k \to \infty} \langle T(u_{n_k}), u_{n_k} \rangle = \langle T(\bar{u}), \bar{u} \rangle$$
    Hemos demostrado que el valor $M$ se alcanza en el vector $\bar{u}$, solo nos falta un detalle crítico: asegurarnos de que $\bar{u}$ mide 1. \\

    \noindent
    Por la convergencia débil, sabemos por propiedad general que $\|\bar{u}\| \le \liminf \|u_{n_k}\| = 1$. ¿Podría ser menor que 1? ¿O incluso cero?
    \begin{itemize}
        \item \tbf{Caso A}: si $M > 0$, si $\|\bar{u}\| < 1$ (y dado que no puede ser cero porque entonces $\langle T\bar{u}, \bar{u} \rangle = 0 \neq M$), podemos definir el vector normalizado $v = \frac{\bar{u}}{\|\bar{u}\|}$, que cumple $\|v\| = 1$. Calculamos su valor:$$\langle Tv, v \rangle = \left\langle T\left(\frac{\bar{u}}{\|\bar{u}\|}\right), \frac{\bar{u}}{\|\bar{u}\|} \right\rangle = \frac{1}{\|\bar{u}\|^2} \langle T\bar{u}, \bar{u} \rangle = \frac{1}{\|\bar{u}\|^2} M$$Como $\|\bar{u}\| < 1$, entonces $\frac{1}{\|\bar{u}\|^2} > 1$. Como $M > 0$, esto implicaría que:$$\langle Tv, v \rangle > M$$¡Contradicción! No podemos encontrar un vector con un valor mayor que el supremo $M$. Por lo tanto, obligatoriamente $\|\bar{u}\| = 1$.
        
        \item \tbf{Caso B}: si $M = 0$Si el supremo es $0$, dado que $T$ es autoadjunto y compacto, por el teorema espectral sus valores propios son reales. Si $M=0$, significa que $\langle Tu, u \rangle \le 0$ para todo $u$, por lo que el operador es semidefinido negativo. En este caso, el valor propio máximo es $0$ y su vector propio asociado (que siempre podemos elegir de norma 1) cumplirá que $\langle T\bar{u}, \bar{u} \rangle = \langle 0, \bar{u} \rangle = 0 = M$.
        
        \item \tbf{Caso C}: si $M < 0$, se razona exactamente igual que en el Caso A, pero la contradicción con las desigualdades se invierte de signo.
    \end{itemize}
    Por lo tanto, en todos los casos existe un vector legítimo con $\|\bar{u}\| = 1$ tal que
    $$M = \max_{\substack{u \in H \\ \|u\|=1}} \langle T(u), u \rangle$$
    Análogamente, repitiendo el proceso con una sucesión minimizante, se demuestra para el ínfimo $m$
    $$m = \min_{\substack{\|u\|=1 \\ u \in H}} \langle T(u), u \rangle$$
\end{ejercicio}


\begin{ejercicio} % //TODO: hacer
    Sea $X = L^2((0,1))$ y $T$ el operador definido por:
    \[ T(u)(x) = \int_{0}^{x} u(t) \, dt, \qquad \forall u \in X \]
    (es el mismo de un ejercicio precedente)

    \begin{itemize}
        \item[i)] Demostrar que $T$ es compacto (\textit{Sugerencia: utilizar adecuadamente el Teorema de Ascoli-Arzelà}).
        
        \item[ii)] ¿El número $0$ pertenece al espectro de $T$? ¿Es un valor propio?
    \end{itemize}
    \tbf{Ejercicio 8.3.2.}
\end{ejercicio}

\begin{ejercicio}
    Sea $X$ un espacio normado, $z \in X$ un elemento suyo fijado y $f$ un funcional en $X^*$. Definimos el siguiente operador $T: X \longrightarrow X$:
    \[ T(x) = f(x)z, \qquad x \in X \]
    Demostrar que $T$ es compacto. \\

    \noindent
    Veámos primero la linealidad y continuidad del operador.
    \begin{itemize}
        \item  \tbf{Linealidad}: tomamos dos elementos $x, y \in X$ y dos escalares $\alpha, \beta \in \mathbb{C}$ (o $\mathbb{R}$). Usando la linealidad del funcional $f \in X^*$
        $$T(\alpha x + \beta y) = f(\alpha x + \beta y)z = \big(\alpha f(x) + \beta f(y)\big)z = \alpha f(x)z + \beta f(y)z = \alpha T(x) + \beta T(y)$$
        Por tanto, $T$ es lineal.
        \item \tbf{Continuidad (= Acotación)}: calculamos la norma de la imagen para cualquier $x \in X$. Como $f(x)$ es un escalar y $z$ es un vector, los escalares salen fuera del módulo de la norma
        $$\|T(x)\| = \|f(x)z\| = |f(x)| \cdot \|z\|$$
        Como el enunciado nos dice que $f \in X^*$, sabemos por definición que $f$ es un funcional continuo, lo que significa que está acotado por su propia norma: $|f(x)| \le \|f\| \cdot \|x\|$, entonces
        $$\|T(x)\| \le \big(\|f\| \cdot \|x\|\big) \cdot \|z\| = \big(\|f\| \cdot \|z\|\big) \cdot \|x\|$$
        y dado que $\|f\|$ y $\|z\|$ son números reales fijos y constantes, si llamamos $C = \|f\| \cdot \|z\|$, tenemos que 
        $$\|T(x)\| \le C \|x\|$$
        por lo que $T$ acotado, y entonces $T$ es continuo.
    \end{itemize}
    Vamos a demostrar ahora que $T$ es \underline{compacto}. Vamos a mirar cuál es la imagen del operador ($\text{Im}(T)$), es decir, qué tipo de vectores pueden salir de aplicar $T$
    $$T(x) = \underbrace{f(x)}_{\text{Escalar}} \cdot \underbrace{z}_{\text{Vector fijo}}$$
    Cualquier vector que salga de aplicar $T$ a un elemento $x$ del espacio es simplemente el vector fijo $z$ multiplicado por un número. Esto significa que toda la imagen de $T$ se encuentra contenida en la recta generada por el vector $z$
    $$\text{Im}(T) = \text{span}\{z\}\implies \text{dim}\big(\text{Im}(T)\big) \le 1 < \infty$$
    por lo que como la imagen tiene dimensión finita, por definición $T$ es un operador de rango finito. Sabemos por un teorema visto que todo operador de rango finito es compacto, pero esto se ve fácilmente de la siguiente manera. \\ \\
    \noindent
    Si tomamos un conjunto acotado $B$ en $X$, entonces su imagen $T(B)$ es un conjunto acotado en $\text{span}\{z\}$, es decir, es un conjunto acotado en una recta. Cualquier conjunto acotado en una recta es relativamente compacto, pues su clausura es compacta (es un segmento de recta cerrado y acotado). Por lo tanto, $T(B)$ es relativamente compacto para cualquier conjunto acotado $B$ de $X$, lo que implica que $T$ es compacto.

\end{ejercicio}
